% ============================================================
% INTRODUCTION GÉNÉRALE
% ============================================================
\chapter*{Introduction générale}
\phantomsection\addcontentsline{toc}{chapter}{Introduction générale}
\markboth{Introduction générale}{Introduction générale}

Un virement de 450\,000~MAD d'une société A vers une société B n'a rien de
suspect. Un virement de B vers C non plus, ni de C vers A. Les trois ensemble,
en quatre minutes, avec une commission de 0,8\,\% prélevée à chaque étape,
forment pourtant la signature d'une \textbf{fraude en carrousel}. Les systèmes
anti-fraude classiques notent chaque transaction ou chaque compte isolément :
par construction, ils ne voient pas une boucle. Pour la voir, il faut
représenter les virements comme un \textbf{graphe} --- les comptes sont des
sommets, les transactions des arêtes --- et chercher des formes dans ce graphe,
pendant que l'argent circule encore.

\begin{tcolorbox}[colback=bleuclair, colframe=bleuprojet, boxrule=0pt, leftrule=3pt,
  arc=0pt, outer arc=0pt, left=8pt, right=8pt, top=6pt, bottom=6pt]
{\sffamily\bfseries\color{bleuprojet}Problématique}\par\smallskip
\emph{Comment détecter, en quelques secondes et sur un flux continu de
virements, les structures de blanchiment qui n'existent qu'à l'échelle du
réseau --- carrousels, cycles longs, comptes mules, dispersions-regroupements,
comptes pivots --- tout en gardant une trace auditable de chaque alerte et
une explication compréhensible par un analyste ?}
\end{tcolorbox}

L'énoncé du Mini-Projet~17 impose la chaîne technique : Kafka pour le flux,
Spark Structured Streaming et GraphFrames pour l'analyse, Delta Lake pour
l'historisation, Streamlit et PyVis pour la restitution. Le travail a consisté
à la faire fonctionner \emph{de bout en bout}, à la \emph{mesurer} contre une
vérité terrain, et à la rendre \emph{rapide} sur un simple ordinateur portable.
La figure~\ref{fig:chiffres} résume ce qui a été obtenu.

\begin{figure}[htp]
\centering
\tuile{bleuprojet}{99,0\,\%}{précision}{101 alertes vraies sur 102}\hfill
\tuile{bleuprojet}{92,9\,\%}{rappel}{39 fraudes détectées sur 42}\hfill
\tuile{orangeprojet}{3,9\,s}{détection d'un carrousel}{médiane, 8~s au maximum}\hfill
\tuile{vertcode}{0}{doublon après panne}{\fichier{kill -9} de Spark en pleine ingestion}
\caption[Les chiffres clés du projet]{Les chiffres clés, mesurés sur le run d'évaluation final
(42~minutes, 50~transactions/s, une fraude toutes les 30~s) et lors du test de panne.}
\label{fig:chiffres}
\end{figure}

Le rapport suit le trajet d'une transaction (figure~\ref{fig:plan}). Le
chapitre~\ref{ch:contexte} cadre le sujet et le chapitre~\ref{ch:etatart}
présente les typologies de blanchiment et les algorithmes de graphe. Le
chapitre~\ref{ch:architecture} décrit l'architecture, le
chapitre~\ref{ch:donnees} les données, le chapitre~\ref{ch:algorithmes} la
détection et le chapitre~\ref{ch:lakehouse} le lakehouse. Le
chapitre~\ref{ch:dashboard} montre le tableau de bord ; les
chapitres~\ref{ch:performance} et~\ref{ch:resultats} mesurent les performances
et la qualité de détection ; les deux derniers traitent des tests, des limites
et du bilan. Presque toutes les figures de données sont tracées à partir des
tables Delta et des journaux de métriques d'un même run réel
(annexe~\ref{ann:reproductibilite}).

\begin{figure}[htp]
\centering
\begin{tikzpicture}[
  etape/.style={rounded corners=3pt, fill=#1, text=white, minimum height=9mm, text width=1.9cm,
    align=center, font=\sffamily\scriptsize\bfseries},
  chap/.style={font=\sffamily\tiny, color=encre2, align=center}]
  \foreach \nom/\ch/\c [count=\i] in {
     {Cadrer\\(sujet, état de l'art)}/{ch.~1--2}/grismoyen,
     {Concevoir\\(architecture)}/{ch.~3}/bleuprojet,
     {Générer\\(données)}/{ch.~4}/bleuprojet,
     {Détecter\\(graphe)}/{ch.~5}/orangeprojet,
     {Historiser\\(Delta Lake)}/{ch.~6}/bleuprojet,
     {Restituer\\(dashboard)}/{ch.~7}/bleuprojet,
     {Mesurer\\(perf., qualité)}/{ch.~8--10}/vertcode} {
    \node[etape=\c] (n\i) at (\i*2.24,0) {\nom};
    \node[chap, below=1mm of n\i] {\ch};
  }
  \foreach \a/\b in {1/2,2/3,3/4,4/5,5/6,6/7}
    \draw[-{Stealth[length=1.6mm]}, line width=0.7pt, color=grismoyen] (n\a) -- (n\b);
\end{tikzpicture}
\caption{Organisation du rapport, calquée sur le trajet d'une transaction.}
\label{fig:plan}
\end{figure}
